\documentclass{article}
\input{../../lib.sty}

\usepackage[backend=biber, style=authoryear-comp, sorting=nyt]{biblatex}
\addbibresource{ref.bib} 

\title{\texttt{fn approximate\_to\_tradeoff}}
\author{\AishwaryaRamasethuAuthor, \YuJuKuAuthor, \JordanAwanAuthor, \MichaelShoemateAuthor}
\begin{document}

\maketitle

\contrib

\section{Hoare Triple}

\subsection*{Precondition}
\subsubsection*{Compiler-verified}
\begin{itemize}
    \item Argument \texttt{epsilon} of type \texttt{f64}
    \item Argument \texttt{delta} of type \texttt{f64}
\end{itemize}

\subsubsection*{Caller-verified}
\texttt{epsilon} is finite and nonnegative, and \texttt{delta} is in $[0,1]$.

\subsection*{Pseudocode}

\lstinputlisting[language=Python,firstline=2,escapechar=`]{./pseudocode/approximate_to_tradeoff.py}

\subsection*{Postcondition}

\begin{theorem}
    The pseudocode returns a function $\tilde{f}$ that is conservative with respect to the intent $I$.
    For \texttt{Reporting}, for all $\alpha \in [0, 1]$,
    \[
        \tilde{f}(\alpha) \le f_{\epsilon,\delta}(\alpha).
    \]
    For \texttt{Calibration}, for all $\alpha \in [0, 1]$,
    \[
        \tilde{f}(\alpha) \ge f_{\epsilon,\delta}(\alpha).
    \]
    Here $f_{\epsilon,\delta}$ is the exact approximate-DP tradeoff curve.
\end{theorem}

We start with the following definition of DP from~\cite{dong2019gaussian}.
\begin{definition}[\protect{\cite[Proposition~2.5]{dong2019gaussian}}]
    \label{def-2-2}
    Let $\epsilon > 0$ and $\delta \geq 0$, and define 
    \begin{equation}
        f_{\epsilon,\delta}(\alpha) = \max(0, 1 - \delta - \exp(\epsilon) \alpha, \exp(-\epsilon) (1 - \delta - \alpha)).
    \end{equation}
    Then we say that a mechanism $M$ satisfies $(\epsilon, \delta)$-DP if it satisfies $f_{\epsilon,\delta}$-DP.
\end{definition}

\begin{proof} 
The definition assumes $\alpha \in [0, 1]$, which is the domain of the returned function.
The code clamps the base of the second affine branch at zero, as in Definition~\ref{def-2-2}.

For \texttt{Reporting}, \texttt{Dir[I]} is upward rounding and \texttt{Opp[I]} is downward rounding.
Thus line \ref{exp-eps} over-estimates $\exp(\epsilon)$, and line \ref{exp-neg-eps} under-estimates $\exp(-\epsilon)$.
The results of these constants are converted exactly into rationals to be used in the tradeoff function.

The function defined on line \ref{tradeoff} implements the formula in Definition~\ref{def-2-2} with exact fractional arithmetic, except that it substitutes these conservative rounded constants for $\exp(\epsilon)$ and $\exp(-\epsilon)$.

For $\alpha \in [0, 1]$, the first branch
\[
    1 - \delta - \texttt{exp\_eps} \cdot \alpha
\]
is no greater than
\[
    1 - \delta - \exp(\epsilon)\alpha,
\]
because $\texttt{exp\_eps} \ge \exp(\epsilon)$ and $\alpha \ge 0$.

For the second branch, if $1 - \delta - \alpha \ge 0$, then
\[
    \texttt{exp\_neg\_eps} \cdot (1 - \delta - \alpha)
    \le
    \exp(-\epsilon)(1 - \delta - \alpha),
\]
because $\texttt{exp\_neg\_eps} \le \exp(-\epsilon)$.
If instead $1 - \delta - \alpha < 0$, then both expressions are negative, so they do not affect the outer maximum with $0$.

Therefore, for all $\alpha \in [0, 1]$, the function returned by the pseudocode is pointwise no greater than $f_{\epsilon,\delta}(\alpha)$.
This is the desired reporting behavior because smaller $\beta$ means a weaker reported tradeoff curve.

For \texttt{Calibration}, the same monotonicity argument reverses.
Now \texttt{Dir[I]} rounds downward and \texttt{Opp[I]} rounds upward, so line \ref{exp-eps} under-estimates $\exp(\epsilon)$ and line \ref{exp-neg-eps} over-estimates $\exp(-\epsilon)$.
The first branch is therefore no smaller than the exact first branch, and the second branch is no smaller than the exact second branch whenever the nonnegative base contributes.
Taking the maximum with zero preserves this pointwise inequality, so $\tilde{f}(\alpha) \ge f_{\epsilon,\delta}(\alpha)$.
This is the desired calibration behavior because larger $\beta$ means a stronger tradeoff curve.
\end{proof}

\printbibliography
\end{document}
